\scanpage{12}
are ``geometric'' fundamental groups --- \unsure{which|which} can moreover be computed by transcendental means, using an embedding of $\overline{\unsure{\deleted{\mathbb Q}K|\deleted{K}\mathbb Q}}$ into $\mathbb C$ (whence an isomorphism $\overline{\mathbb Q}\simeq\overline{\mathbb Q}_0$), as the profinite completions of the $\pi_1(U_i(\underset{\unsure{\mathbb Q|k|\mathbb C}}{\mathbb C}),\bar\eta)$, \unsure{in|where} now $\bar\eta$ is interpreted as a \deleted{\unsure{complex|complex}} point common to the topological spaces underlying the complex varieties $U_i(\mathbb C)$.

The exact sequence (8) is therefore the inverse limit of the homotopy exact sequences (12), which gives in particular
\[
\begin{aligned}
(13)\qquad \pi_{\bar K/K}&\simeq\varprojlim_i\underbrace{\pi_1(\bar U_i,\bar\eta)}_{\pi_i}\\
&\simeq\varprojlim_i\pi_1(U_i(\mathbb C),\bar\eta)^{\wedge}
\end{aligned}
\]
\margin{This interpretation is obviously not particular to the base field $\mathbb Q$ --- it holds \unsure{for|for} every field extension $K/K_0$}

Using the fibrations of the $U_i$ (in the case where one restricts oneself to taking Artin's elementary varieties), one finds that every $\pi_i$ is a group \deleted{\unsure{by|pr}} successive extension of \inserted{$n$} \uline{free} profinite groups (\unsure{here|where} $n=\dim U_i=\operatorname{deg.tr.}K/\mathbb Q$). This gives again
\margin{11}
